\documentclass[12pt]{article}
\usepackage[T1]{fontenc}
\usepackage[utf8]{inputenc}
\usepackage{lmodern}
\usepackage{textcomp}
\usepackage{enumitem}
\usepackage{geometry}
\geometry{margin=1in}
\begin{document}
\begin{center}
Answer Key - History - Undergraduate Level Assessment 6 \
\small (Answers are ordered exactly as in the question paper.)
\end{center}

\section*{Multiple Choice}
\begin{enumerate}[label=\arabic*.]
\item \textbf{Answer:} D
\\ \textit{Options:} A) to demonstrate their military prowess to potential allies and enemies.; B) to expand their empire's territory.; C) to unite their own population through a common enemy.; D) to identify and cultivate talented leaders.; E) to gain control of vital resources.
\\ \textit{Note:} The Inca engaged in warfare not only to expand their territory but also as a means to identify and promote skilled individuals into leadership roles, thereby strengthening their societal and military structure.
\item \textbf{Answer:} E
\\ \textit{Options:} A) Hohokam; B) Folsom; C) Adena; D) Ancestral Puebloan; E) Mississippian
\\ \textit{Note:} The Mississippian culture, recognized for their construction of large temple mounds, relied heavily on maize and squash agriculture as key staples in their diet, which facilitated their complex societal structures and trade networks.
\item \textbf{Answer:} A
\\ \textit{Options:} A) Mississippi; B) Rocky Mountain; C) Appalachian; D) McKenzie Corridor; E) Pacifica
\\ \textit{Note:} The correct answer is Mississippi because it refers to the extensive glacial deposits and the influence of the Laurentide Ice Sheet during the Pleistocene era, which affected the geography and hydrology of the region, including the Mississippi River's formation.
\item \textbf{Answer:} A
\\ \textit{Options:} A) pelagic; lacustrine; B) pelagic; littoral; C) pelagic; midden; D) littoral; lacustrine; E) midden; littoral
\\ \textit{Note:} The correct answer is based on the understanding that the Lake Forest Archaic tradition primarily utilized resources from the open sea, referred to as pelagic, while the Maritime Archaic focused on hunting freshwater creatures, known as lacustrine, found in lakes.
\item \textbf{Answer:} B
\\ \textit{Options:} A) the presence of a nearby museum or storage facility.; B) naturally exposed layers in which to search for artifacts, ecofacts, or bone.; C) the accessibility of the site by public transportation.; D) the presence of primary and secondary refuse.; E) the existence of a cache or collection of artifacts.
\\ \textit{Note:} The appropriateness and usefulness of pedestrian survey is primarily determined by the presence of naturally exposed layers, as these layers facilitate the visibility and accessibility of artifacts, ecofacts, or bone, making it easier for archaeologists to conduct effective surveys and gather relevant historical data.
\end{enumerate}

\section*{True / False}
\begin{enumerate}[label=\arabic*.]
\setcounter{enumi}{5}
\item \textbf{Answer:} False
\\ \textit{Options:} A) True; B) False
\\ \textit{Note:} The correct answer is: the Yangtze River and mountains
\item \textbf{Answer:} True
\\ \textit{Options:} A) True; B) False
\\ \textit{Note:} According to the passage: two
\item \textbf{Answer:} True
\\ \textit{Options:} A) True; B) False
\\ \textit{Note:} According to the passage: GamePro and EGM
\item \textbf{Answer:} True
\\ \textit{Options:} A) True; B) False
\\ \textit{Note:} According to the passage: eight local factories and institutions
\item \textbf{Answer:} False
\\ \textit{Options:} A) True; B) False
\\ \textit{Note:} The correct answer is: Alfonso Soriano
\end{enumerate}

\section*{Long Form Response}
\begin{enumerate}[label=\arabic*.]
\setcounter{enumi}{10}
\item \textbf{Answer:} iPod Touch and iPhone
\\ \textit{Note:} Expected answer: iPod Touch and iPhone
\item \textbf{Answer:} Palestine
\\ \textit{Note:} Expected answer: Palestine
\end{enumerate}

\vfill
\noindent\textit{End of Answer Key}
\end{document}